In this chapter we study the \textbf{infinite-horizon linear quadratic regulator} from a data-driven
viewpoint. The starting point is the standard stochastic LQR problem for the discrete-time
linear system
\[
x_{t+1} = A x_t + B u_t,
\qquad x_0 \sim \mathcal D,
\]
where the objective is to minimize the expected infinite-horizon quadratic cost
\[
\min_{\{u_t\}} \;
\mathbb E\!\left[\sum_{t=0}^{\infty}
\left(x_t^\top Q x_t + u_t^\top R u_t\right)\right].
\]
Throughout, we assume
\[
Q \succeq 0,
\qquad R \succ 0,
\]
and we require the standard structural assumptions
\[
(A,B)\ \text{stabilizable},
\qquad
(Q^{1/2},A)\ \text{detectable}.
\]
Under these assumptions, the infinite-horizon LQR problem is well posed and admits an
\textbf{optimal stationary policy} of the form
\[
u_t = K^\star x_t,
\]
that is, the optimal controller is a linear state-feedback law.\\

When the matrices $A$ and $B$ are known, the optimal gain $K^\star$ can be computed in
different but equivalent ways. A first route is dynamic programming, which leads to the
Riccati recursion and, in the infinite-horizon limit, to the algebraic Riccati equation.
A second route is therefore to solve directly the algebraic Riccati equation and recover the
optimal gain from its stabilizing solution. A third route is based on a \textbf{semidefinite programming
reformulation}. These viewpoints are equivalent at the model-based level, but they suggest
different ways of reformulating the problem when the model is not available exactly and one
would like to reason directly from data.

\subsection{LQR parameterization}

A useful perspective for data-driven control is to parameterize the LQR problem directly in
terms of the feedback matrix $K$. Instead of optimizing over arbitrary input sequences, we
restrict attention to linear state-feedback policies
\[
u_t = K x_t,
\]
and \textbf{optimize over all stabilizing gains} $K$.\\
To make the dependence on $K$ explicit, let us assume for simplicity that the initial condition
satisfies
\[
\mathbb E[x_0 x_0^\top] = I.
\]
If the feedback law $u_t = Kx_t$ is applied, then the closed-loop system becomes
\[
x_{t+1} = (A+BK)x_t.
\]
For a stabilizing gain $K$, it is natural to introduce the infinite-horizon state covariance
\[
P \;=\; \mathbb E\!\left[\sum_{t=0}^{\infty} x_t x_t^\top\right].
\]
This matrix collects the second-order state information accumulated over the whole trajectory.
Using the closed-loop dynamics, one obtains
\[
P
= \mathbb E[x_0x_0^\top]
+ (A+BK)\,\mathbb E[x_0x_0^\top]\,(A+BK)^\top
+ \cdots,
\]
and therefore, since $\mathbb E[x_0x_0^\top]=I$, the matrix $P$ \textbf{satisfies the Lyapunov equation}
\[
P = I + (A+BK)P(A+BK)^\top.
\]
This identity is fundamental. In particular, it shows that $P$ is the unique positive definite
solution of the Lyapunov equation \textbf{if and only if} the closed-loop matrix $A+BK$ is Schur
stable.\\
Using this parameterization, the infinite-horizon LQR cost can be rewritten in terms of $K$
and $P$. Indeed, since $u_t = Kx_t$, we have
\[
x_t^\top Q x_t + u_t^\top R u_t
= x_t^\top (Q + K^\top R K)x_t,
\]
hence summing over time and taking expectations gives
\[
J(K)
=
\mathrm{Tr}(QP) + \mathrm{Tr}(K^\top R K\, P).
\]
Therefore, the \textbf{infinite-horizon LQR problem can be expressed as}
\[
\begin{aligned}
\min_{K,\;P \succ 0}\quad
& \mathrm{Tr}(QP) + \mathrm{Tr}(K^\top R K\,P) \\
\text{s.t.}\quad
& P = I + (A+BK)P(A+BK)^\top .
\end{aligned}
\]
This formulation is important because it makes the optimization variable $K$ explicit and
\textbf{replaces the dynamical system by a matrix equality constraint}. At the same time, it also makes clear that the problem is \emph{non-convex}: the Lyapunov constraint is bilinear in $K$ and $P$, and the objective also contains the product of these variables.\\
This non-convex parameterization is nevertheless extremely useful, because it is the form from
which many data-driven LQR methods start. The \textbf{main challenge} of data-driven LQR is precisely to recover or approximate the optimal gain $K^\star$ \textbf{without explicitly knowing the matrices} $A$ and $B$, while preserving the stabilizing and performance properties of the model-based solution.

\subsection{SDP formulation of the LQR}

The parameterization introduced above is very convenient conceptually, but it still leads to a
\emph{non-convex} optimization problem:
\[
\begin{aligned}
\min_{K,\;P \succ 0}\quad
& \text{Tr}(QP) + \text{Tr}(K^\top R K\,P) \\
\text{s.t.}\quad
& P = I + (A+BK)P(A+BK)^\top .
\end{aligned}
\]
The non-convexity appears in two places. First, the Lyapunov equality constraint is bilinear in
$K$ and $P$. Second, the term $\text{Tr}(K^\top R K\,P)$ also contains products of the decision
variables. A remarkable fact is that this LQR problem can nevertheless be reformulated
\emph{exactly} as a \textbf{convex semidefinite program} (SDP).\\
Before deriving the SDP formulation, it is useful to recall a standard tool that is repeatedly used
in control and convex optimization.

\paragraph{Schur complement.}
Consider a symmetric block matrix
\[
X=\begin{bmatrix}
A & B\\
B^\top & C
\end{bmatrix}.
\]
If $C$ is invertible, the matrix
\[
A-BC^{-1}B^\top
\]
is called the \textbf{Schur complement} of $C$ in $X$.
The corresponding lemma states that, if $C\succ 0$, then
\[
\begin{bmatrix}
A & B\\
B^\top & C
\end{bmatrix}\succeq 0
\qquad \Longleftrightarrow \qquad
A-BC^{-1}B^\top \succeq 0.
\]
This result allows us to transform matrix inequalities involving inverses into linear matrix
inequalities, which are exactly the constraints used in semidefinite programming.

\medskip

We now \textbf{derive the SDP formulation} of LQR step by step. Start from
\begin{equation}
\begin{aligned}
\min_{K,\;P \succ 0}\quad
& \text{Tr}(QP) + \text{Tr}(K^\top R K\,P) \\
\text{s.t.}\quad
& P = I + (A+BK)P(A+BK)^\top .
\end{aligned}
\label{eq:original}
\end{equation}
A first observation is that we can \textbf{replace the equality constraint by the inequality}
\[
P \succeq I + (A+BK)P(A+BK)^\top
\]
without changing the optimum. This gives the equivalent problem
\begin{equation}
\begin{aligned}
\min_{K,\;P \succ 0}\quad
& \text{Tr}(QP) + \text{Tr}(K^\top R K\,P) \\
\text{s.t.}\quad
& P \succeq I + (A+BK)P(A+BK)^\top .
\end{aligned}
\label{eq:relaxed}
\end{equation}
At first sight, this equivalence may look surprising, so it is worth understanding why it is true.\\
Fix a feedback gain $K$ such that $A+BK$ is Schur stable, and let $P_K$ denote the unique
positive definite solution of the Lyapunov equation
\begin{equation}
P_K = I + (A+BK)P_K(A+BK)^\top.
\label{eq:Lyapunov}
\end{equation}
Now consider any matrix $P$ satisfying the inequality
\begin{equation}
P \succeq I + (A+BK)P(A+BK)^\top.
\label{eq:inequality}
\end{equation}
Define the difference
\[
Y := P - P_K.
\]
Subtracting \eqref{eq:Lyapunov} from \eqref{eq:inequality}, we obtain
\[
Y \succeq (A+BK)\,Y\,(A+BK)^\top.
\]
Iterating this inequality yields
\[
Y \succeq (A+BK)\,Y\,(A+BK)^\top
\succeq (A+BK)^2\,Y\,\big((A+BK)^\top\big)^2
\succeq \cdots \succeq 0.
\]
Hence $Y \succeq 0$, that is,
\[
P \succeq P_K.
\]
Therefore, for a fixed stabilizing $K$, the smallest feasible matrix $P$ is exactly the Lyapunov
solution $P_K$. Since the objective is linear and increasing in $P$ (recall that $Q\succeq 0$ and
$R\succ 0$), the optimizer will always choose the smallest feasible $P$, so the inequality in
\eqref{eq:relaxed} is active at the optimum. This explains why \eqref{eq:original} and \eqref{eq:relaxed} are equivalent.

\medskip

The next step is to \textbf{eliminate the bilinear term in the constraint} by introducing the change of
variables
\[
L := KP.
\]
Using this substitution, the constraint in \eqref{eq:relaxed} becomes
\[
P - I - (AP+BL)P^{-1}(AP+BL)^\top \succeq 0.
\]
This expression is still nonlinear because of the inverse $P^{-1}$, but now it is in a form where
the Schur complement can be applied. Since $P\succ 0$, the Schur complement lemma implies
that the inequality above is equivalent to the linear matrix inequality
\begin{equation}
\begin{bmatrix}
P-I & AP+BL\\
(AP+BL)^\top & P
\end{bmatrix}
\succeq 0.
\end{equation}
So the Lyapunov inequality has now been converted into an Linear Matrix Inequality, which is convex in the
variables $P$ and $L$.\\

We now turn to the objective function. Since $L=KP$, one would like to replace the term
$KPK^\top$ by a new matrix variable. Introduce
\[
M := KPK^\top.
\]
Then
\[
\text{Tr}(K^\top RKP)=\text{Tr}(RKPK^\top)=\text{Tr}(RM),
\]
and the objective becomes
\[
\text{Tr}(QP)+\text{Tr}(RM).
\]
However, the identity $M=KPK^\top$ is again nonlinear. As before, we relax it to the inequality
\[
M \succeq KPK^\top.
\]
This leads to
\[
M - KPK^\top \succeq 0.
\]
Using $L=KP$, this can be rewritten as
\[
M - LP^{-1}L^\top \succeq 0.
\]
Applying again the Schur complement with respect to $P\succ 0$, we obtain the equivalent LMI
\begin{equation}
\begin{bmatrix}
M & L\\
L^\top & P
\end{bmatrix}
\succeq 0.
\label{eq:Schur}
\end{equation}
As in the previous step, this relaxation is tight at the optimum because the objective is linear in
$M$ and $R\succ 0$, so the optimizer will choose the smallest feasible $M$.\\
Putting everything together, we arrive at the \textbf{SDP formulation of the infinite-horizon LQR}:
\begin{equation}
\begin{aligned}
\min_{L,M,P}\quad
& \text{Tr}(QP)+\text{Tr}(RM)\\
\text{s.t.}\quad
& \begin{bmatrix}
M & L\\
L^\top & P
\end{bmatrix}\succeq 0,\\[1em]
& \begin{bmatrix}
P-I & AP+BL\\
(AP+BL)^\top & P
\end{bmatrix}\succeq 0,\\
& P \succ 0.
\end{aligned}
\end{equation}
This is now a \textbf{convex optimization problem}: the objective is affine in the decision variables
and the constraints are positive semidefinite constraints, that is, linear matrix inequalities.\\

Once the problem has been solved, the optimal feedback gain is recovered from
\[
K = LP^{-1}.
\]
So, although the original LQR problem was written in a non-convex form in terms of $K$ and
$P$, it admits an equivalent convex reformulation in the lifted variables $(L,M,P)$.\\

This result is much more than a technical curiosity. Semidefinite programming plays a central
role in modern control because many synthesis and analysis problems can be written in terms of
LMIs after a suitable change of variables. The LQR problem is one of the cleanest examples of
this idea: a nonlinear controller synthesis problem can be converted into a convex optimization
problem for which global optimality is guaranteed.\\
For us, this SDP viewpoint is particularly useful because it opens the door to \textbf{data-driven}
reformulations. Indeed, once the model-based LQR problem has been written in terms of matrix
inequalities, one can try to replace the quantities involving $A$ and $B$ by quantities constructed
directly from data.

\subsection{From model-based to data-driven LQR}

At this point it is useful to step back and summarize the situation. For the infinite-horizon
LQR problem, we have seen three equivalent \emph{model-based} solution routes:
\begin{itemize}
    \item dynamic programming;
    \item the algebraic Riccati equation;
    \item the SDP formulation.
\end{itemize}
Although these approaches look quite different, they all share the same fundamental
requirement: \textbf{they need the matrices $A$ and $B$ to be known}. Dynamic programming
needs the model to write the Bellman recursion, the Riccati approach needs the model to form
the algebraic Riccati equation, and the SDP formulation also contains $A$ and $B$ explicitly in
its constraints. Therefore, if the system matrices are not available, none of these procedures can
be applied directly.\\
This observation motivates the central question of this chapter:
\begin{quote}
\emph{Can we design the optimal LQR controller directly from data, without first having an
exact model of the system?}
\end{quote}
This question belongs to the broader area of \textbf{data-driven control}. At a high level, one can
distinguish two main philosophies.

\paragraph{Indirect data-driven control.}
In the indirect approach, one first uses data to identify a model of the system, together with a
description of the uncertainty, and then one designs a controller based on the identified model.
In short, the logic is
\[
\text{data} \;\longrightarrow\; \text{model identification} + \text{uncertainty quantification}
\;\longrightarrow\; \text{control design}.
\]
This is the classical modular viewpoint: first system identification, then control synthesis.

\paragraph{Direct data-driven control.}
A different philosophy is to \emph{bypass explicit model identification} and attempt to compute a
controller directly from data. This is the so-called \textbf{direct} data-driven approach. In recent
years this viewpoint has attracted a lot of attention, partly because it promises to avoid the
intermediate identification step and to exploit data more directly for control design.\\

However, this promise comes with several trade-offs. In particular, when comparing indirect and
direct approaches, one should keep in mind issues such as:
\begin{itemize}
    \item modularity versus non-modularity;
    \item computational tractability;
    \item optimality versus suboptimality;
    \item robustness and robust stability guarantees;
    \item the amount of data required.
\end{itemize}
For the LQR problem, these questions can be addressed quite explicitly, which is one of the
reasons why LQR is an ideal benchmark for studying data-driven control design.

\subsection{Indirect data-driven control}
\subsubsection{System identification from input-state data}

Let us start from the indirect route. Suppose that the system evolves according to
\[
x_{t+1} = A x_t + B u_t + w_t,
\]
where $w_t$ is a process disturbance or modeling error. Assume that we collect a trajectory of
state-input data and arrange it into the matrices
\[
X_0 = \begin{bmatrix} x_0 & x_1 & \cdots & x_{T-1} \end{bmatrix},
\qquad
U_0 = \begin{bmatrix} u_0 & u_1 & \cdots & u_{T-1} \end{bmatrix},
\]
\[
W_0 = \begin{bmatrix} w_0 & w_1 & \cdots & w_{T-1} \end{bmatrix},
\qquad
X_1 = \begin{bmatrix} x_1 & x_2 & \cdots & x_T \end{bmatrix}.
\]
By stacking the dynamics over the whole trajectory, we obtain the compact relation
\[
X_1 = A X_0 + B U_0 + W_0.
\]
This matrix equation is the starting point for identification from data. It tells us that the
unknown matrices $A$ and $B$ must explain, as well as possible, the one-step transitions
observed in the data.\\
A fundamental requirement is that the data be sufficiently informative. In this context, the
relevant notion is \textbf{persistence of excitation}. More precisely, the data are said to be
sufficiently exciting if
\[
\rank\!\begin{bmatrix}
U_0\\[0.3em]
X_0
\end{bmatrix}
= n+m,
\]
where $n$ is the state dimension and $m$ is the input dimension.
This rank condition means that the collected trajectory excites all state and input directions
needed to identify the dynamics. If the data do not satisfy this condition, then different pairs
$(A,B)$ may be consistent with the same observations, and the model cannot be uniquely
reconstructed.\\

Given the data, a natural identification procedure is ordinary least squares. One solves
\[
[\hat B,\hat A]
=
\arg\min_{B,A}
\left\|
X_1 - [B,\;A]
\begin{bmatrix}
U_0\\[0.3em]
X_0
\end{bmatrix}
\right\|_F.
\]
Equivalently, one searches for the matrices $A$ and $B$ that best fit the observed one-step
state transitions in the Frobenius norm sense. Under the persistence of excitation condition, this
least-squares problem has a unique solution. Hence, with sufficiently rich data, one can recover
an estimate $(\hat A,\hat B)$ of the system matrices.

\subsubsection{Indirect and certainty-equivalent LQR}

Once the model has been identified, the most immediate strategy is simply to replace the true
unknown matrices $(A,B)$ by their estimates $(\hat A,\hat B)$ and then solve the LQR problem as
if the estimated model were exact. This is the classical \textbf{certainty-equivalence principle}.\\
Using the LQR parameterization introduced earlier, the certainty-equivalent design solves
\[
\begin{aligned}
\min_{K,\;P \succ 0}\quad
& \text{Tr}(QP) + \text{Tr}(K^\top R K\,P) \\
\text{s.t.}\quad
& P = I + (\hat A+\hat B K)P(\hat A+\hat B K)^\top,
\end{aligned}
\]
where the identified model is obtained from the least-squares problem
\[
[\hat B,\hat A]
=
\arg\min_{B,A}
\left\|
X_1 - [B,\;A]
\begin{bmatrix}
U_0\\[0.3em]
X_0
\end{bmatrix}
\right\|_F.
\]
In words, the procedure is:
\begin{enumerate}
    \item collect input-state data;
    \item identify $(\hat A,\hat B)$ from the data;
    \item design the LQR controller for the estimated model.
\end{enumerate}
This is why the method is called \textbf{indirect}: the data are used only to build a model, and
the controller is then synthesized by a standard model-based technique.\\
The certainty-equivalent approach is widely used both in theory and in applications. The main
reasons are clear:
\begin{itemize}
    \item it is \textbf{modular}, because identification and control design are separated;
    \item it is \textbf{tractable}, since after identification one can use the Riccati equation or the SDP
    formulation developed earlier;
    \item it requires only the minimum amount of excitation needed to identify the model, namely
    enough data to recover $n+m$ independent directions.
\end{itemize}
At the same time, it also has \textbf{clear limitations}. Since the controller is designed for the estimated
model rather than the true one, the resulting gain is generally only \textbf{suboptimal} for the real
system. Moreover,\textbf{ robustness} is only moderate unless one explicitly accounts for identification
errors and uncertainty in the design step. In other words, certainty-equivalence is simple and
powerful, but it does not by itself provide a complete answer to the issue of robustness.

This discussion highlights a key tension in data-driven control. The indirect route is attractive
because it preserves the classical architecture of \emph{identification first, control second}, and
therefore remains easy to analyze and implement. On the other hand, it may fail to fully exploit
the information in the data for the actual control objective.

\subsection{Direct data-driven LQR via trajectory parameterization}

The indirect approach discussed above follows the classical two-step pipeline
\[
\text{data} \;\longrightarrow\; (\hat A,\hat B) \;\longrightarrow\; K.
\]
A natural question is whether one can avoid the explicit identification step and parameterize the
controller directly in terms of the collected data. This leads to a \textbf{direct} data-driven
formulation of LQR.\\
The key idea is to use the input-state trajectory itself as a parameterization device. Recall that,
for the noisy system
\[
x_{t+1} = A x_t + B u_t + w_t,
\]
the stacked data matrices satisfy
\[
X_1 = A X_0 + B U_0 + W_0.
\]
Assume again that the data are persistently exciting, namely
\[
\rank\!\begin{bmatrix}
U_0\\[0.3em]
X_0
\end{bmatrix}=n+m.
\]
This rank condition has an important consequence: for \emph{every} feedback gain $K$ there
exists a matrix $G$ such that
\begin{equation}
\begin{bmatrix}
K\\[0.3em]
I
\end{bmatrix}
=
\begin{bmatrix}
U_0\\[0.3em]
X_0
\end{bmatrix}G.
\label{eq:KGparam}
\end{equation}
So the pair $(K,I)$ can be represented as a linear combination of the columns of the measured
data matrix. This is the basic trajectory parameterization.\\
Once \eqref{eq:KGparam} is available, the closed-loop matrix can also be written directly in terms
of the data. Indeed,
\[
A+BK
=
[B\;\;A]
\begin{bmatrix}
K\\[0.3em]
I
\end{bmatrix}
=
[B\;\;A]
\begin{bmatrix}
U_0\\[0.3em]
X_0
\end{bmatrix}G.
\]
Using the data equation
\[
[B\;\;A]
\begin{bmatrix}
U_0\\[0.3em]
X_0
\end{bmatrix}
=
X_1-W_0,
\]
we obtain
\begin{equation}
A+BK = (X_1-W_0)G.
\label{eq:closedloopdata}
\end{equation}
This identity is crucial: it shows that the unknown closed-loop matrix can be represented from
data through the variable $G$.\\
If the noise term $W_0$ is neglected, the relation simplifies to
\[
A+BK = X_1G.
\]
This is exactly the certainty-equivalent viewpoint in direct form: \textbf{one behaves as if the measured
data were generated by a noise-free system}. In that sense, the indirect certainty-equivalent
approach and the direct certainty-equivalent approach are closely related.\\

Substituting the relation $A+BK=X_1G$ into the LQR parameterization yields a direct
data-driven reformulation:
\[
\begin{aligned}
\min_{P \succeq I,\;K,\;G}\quad
& \text{Tr}(QP)+\text{Tr}(K^\top RKP)\\
\text{s.t.}\quad
& X_1GPG^\top X_1^\top - P + I \preceq 0,\\
& \begin{bmatrix}
K\\[0.3em]
I
\end{bmatrix}
=
\begin{bmatrix}
U_0\\[0.3em]
X_0
\end{bmatrix}G.
\end{aligned}
\]
The first constraint is simply the Lyapunov inequality written using the data-based expression
for the closed-loop matrix, while the second constraint imposes consistency between the gain
$K$ and the trajectory parameterization $G$.\\
A useful \textbf{interpretation} is obtained by solving the second constraint explicitly. Since
\[
\begin{bmatrix}
K\\[0.3em]
I
\end{bmatrix}
=
\begin{bmatrix}
U_0\\[0.3em]
X_0
\end{bmatrix}G,
\]
we get
\[
K = U_0G,
\qquad
I = X_0G.
\]
Hence the optimization variables can be reduced to $(G,P)$ and the problem becomes
\[
\begin{aligned}
\min_{G,\;P \succ 0}\quad
& \text{Tr}(QP)+\text{Tr}\!\big(G^\top U_0^\top R U_0 G P\big)\\
\text{s.t.}\quad
& X_1GPG^\top X_1^\top - P + I \preceq 0,\\
& X_0G = I.
\end{aligned}
\]
This makes clear that $G$ itself plays the role of the \textbf{policy parameter}. Indeed, the control law is
\[
u_t = Kx_t = U_0Gx_t,
\]
and, correspondingly, the closed-loop dynamics become
\[
x_{t+1} = X_1Gx_t,
\]
again in the noise-free certainty-equivalent interpretation.

\subsubsection{Non-uniqueness of the trajectory parameterization}

An important feature of this direct formulation is that the matrix $G$ is generally \emph{not
unique}. The reason is that the linear system
\[
\begin{bmatrix}
K\\[0.3em]
I
\end{bmatrix}
=
\begin{bmatrix}
U_0\\[0.3em]
X_0
\end{bmatrix}G
\]
may have infinitely many solutions whenever the data matrix
\[
\begin{bmatrix}
U_0\\[0.3em]
X_0
\end{bmatrix}
\]
has more columns than rows. In that case, the set of feasible $G$ contains an affine subspace
whose direction is precisely the nullspace
\[
\ker\!\begin{bmatrix}
U_0\\[0.3em]
X_0
\end{bmatrix}.
\]
So two different matrices $G$ may generate the same gain $K$ and still satisfy the consistency
constraint.\\
This non-uniqueness is not just a technical detail. It means that the direct formulation admits
extra degrees of freedom that are invisible at the level of $K$. Those degrees of freedom can be
used either in a favorable way or in a harmful way, depending on how the optimization problem
is regularized. This is why, in practice, one usually augments the direct formulation with a
\textbf{regularizer} that selects one particular representative among all feasible $G$.\\

Broadly speaking, two types of regularization are of special interest:
\begin{itemize}
    \item a \textbf{certainty-equivalence-promoting} regularizer, which selects the solution most closely
    related to the least-squares model estimate;
    \item a \textbf{robustness-promoting} regularizer, which exploits the extra freedom in $G$ to improve
    robustness with respect to noise and uncertainty.
\end{itemize}
So, in the direct approach, regularization is not only a numerical device: it is part of the control
design itself.

\subsubsection{Equivalence between direct and indirect certainty-equivalent LQR}

The previous discussion suggests that there should be a precise connection between the direct
and indirect approaches. In fact, under a suitable choice of regularization, the direct
certainty-equivalent formulation is equivalent to the indirect one based on least-squares
identification.\\

To see the mechanism, recall that among all solutions of
\[
\begin{bmatrix}
K\\[0.3em]
I
\end{bmatrix}
=
\begin{bmatrix}
U_0\\[0.3em]
X_0
\end{bmatrix}G,
\]
one can single out the minimum-norm solution
\[
G =
\begin{bmatrix}
U_0\\[0.3em]
X_0
\end{bmatrix}^{\dagger}
\begin{bmatrix}
K\\[0.3em]
I
\end{bmatrix},
\]
where ${}^\dagger$ denotes the pseudoinverse. Equivalently, one can enforce
that $G$ be orthogonal to the nullspace of the data matrix. Writing
\[
\Pi := I -
\begin{bmatrix}
U_0\\[0.3em]
X_0
\end{bmatrix}^{\dagger}
\begin{bmatrix}
U_0\\[0.3em]
X_0
\end{bmatrix},
\]
the matrix $\Pi$ is the orthogonal projector onto
\[
\ker\!\begin{bmatrix}
U_0\\[0.3em]
X_0
\end{bmatrix}.
\]
Therefore, the orthogonality condition can be written as
\[
\Pi G = 0.
\]
It is useful to remember the main properties of this projector:
\[
\Pi^2 = \Pi,
\qquad
\Pi^\top = \Pi,
\qquad
\begin{bmatrix}
U_0\\[0.3em]
X_0
\end{bmatrix}\Pi = 0.
\]
Now substitute the minimum-norm expression for $G$ into the direct formulation. Since
\[
G =
\begin{bmatrix}
U_0\\[0.3em]
X_0
\end{bmatrix}^{\dagger}
\begin{bmatrix}
K\\[0.3em]
I
\end{bmatrix},
\]
we obtain
\[
X_1G
=
X_1
\begin{bmatrix}
U_0\\[0.3em]
X_0
\end{bmatrix}^{\dagger}
\begin{bmatrix}
K\\[0.3em]
I
\end{bmatrix}.
\]
On the other hand, the least-squares estimates introduced before satisfy
\[
[\hat B\;\;\hat A]
=
X_1
\begin{bmatrix}
U_0\\[0.3em]
X_0
\end{bmatrix}^{\dagger}.
\]
Hence
\[
X_1G
=
[\hat B\;\;\hat A]
\begin{bmatrix}
K\\[0.3em]
I
\end{bmatrix}
=
\hat A+\hat B K.
\]
This shows that the data-based closed-loop matrix appearing in the direct approach coincides
exactly with the identified closed-loop matrix of the indirect least-squares model.

As a consequence, once the orthogonality constraint $\Pi G=0$ is imposed, the direct
certainty-equivalent problem becomes equivalent to the indirect certainty-equivalent LQR problem
\[
\begin{aligned}
\min_{P \succeq I,\;K}\quad
& \text{Tr}(QP)+\text{Tr}(K^\top RKP)\\
\text{s.t.}\quad
& (\hat A+\hat B K)P(\hat A+\hat B K)^\top - P + I \preceq 0,
\end{aligned}
\]
with
\[
[\hat B\;\;\hat A]
=
\arg\min_{B,A}
\left\|
X_1 - [B\;\;A]
\begin{bmatrix}
U_0\\[0.3em]
X_0
\end{bmatrix}
\right\|_F.
\]
So the indirect least-squares route is not disconnected from the direct one: it corresponds to a
specific way of choosing the trajectory parameterization, namely the \textbf{minimum-norm choice of
$G$ orthogonal to the nullspace of the data matrix}.\\
This equivalence is conceptually important. It tells us that certainty-equivalence is not tied to
the explicit identification of a model; rather, it can also be understood as a particular
regularization of the direct data-driven parameterization. At the same time, it clarifies where
direct approaches may depart from indirect ones: by exploiting the nullspace freedom in $G$,
one may design regularizers that do \emph{not} reproduce least-squares identification and that
instead aim at improved robustness.

\subsubsection{Certainty-equivalence-promoting regularization}

We now make the previous equivalence more operational. In the trajectory-parameterized direct
formulation, the ambiguity in $G$ comes from the nullspace of the data matrix
\[
\begin{bmatrix}
U_0\\[0.3em]
X_0
\end{bmatrix}.
\]
If one wants the direct solution to reproduce the same behavior as the indirect least-squares
approach, then one should select the representative of $G$ that is orthogonal to this nullspace.\\
Recall the projector
\[
\Pi := I -
\begin{bmatrix}
U_0\\[0.3em]
X_0
\end{bmatrix}^{\dagger}
\begin{bmatrix}
U_0\\[0.3em]
X_0
\end{bmatrix},
\]
which is the orthogonal projector onto
\[
\ker\!\begin{bmatrix}
U_0\\[0.3em]
X_0
\end{bmatrix}.
\]
The exact orthogonality constraint
\[
\Pi G = 0
\]
forces $G$ to belong to the orthogonal complement of the nullspace, and therefore selects the
minimum-norm solution
\[
G =
\begin{bmatrix}
U_0\\[0.3em]
X_0
\end{bmatrix}^{\dagger}
\begin{bmatrix}
K\\[0.3em]
I
\end{bmatrix}.
\]
As we have seen, this is precisely the choice that makes the direct certainty-equivalent
formulation equivalent to the indirect least-squares LQR design.\\
Instead of imposing the hard constraint $\Pi G=0$, one can \textbf{lift it to the objective} and use
it as a regularizer. This leads to the optimization problem
\[
\begin{aligned}
\min_{P \succeq I,\;K,\;G}\quad
& \text{Tr}(QP)+\text{Tr}(K^\top RKP)+\lambda \|\Pi G\| \\
\text{s.t.}\quad
& X_1GPG^\top X_1^\top - P + I \preceq 0,\\
& \begin{bmatrix}
K\\[0.3em]
I
\end{bmatrix}
=
\begin{bmatrix}
U_0\\[0.3em]
X_0
\end{bmatrix}G.
\end{aligned}
\]
Here $\lambda \ge 0$ is a tuning parameter. When $\lambda$ is large, the optimizer is strongly
encouraged to make $\Pi G$ small, that is, to choose $G$ almost orthogonal to the nullspace of
the data matrix. In the limit, this reproduces the same effect as the hard orthogonality
constraint.\\

This regularizer is therefore \textbf{certainty-equivalence-promoting}: it biases the direct
data-driven problem toward the same solution that would be obtained by first identifying a
least-squares model and then designing the LQR controller for that model.

The parameter $\lambda$ has a clear interpretation. It interpolates between two extremes:
\begin{itemize}
    \item if $\lambda=0$, the optimizer is free to exploit the whole family of feasible matrices $G$,
    including the nullspace directions;
    \item if $\lambda$ is very large, the optimizer is pushed toward the minimum-norm solution
    associated with least-squares system identification.
\end{itemize}
In this sense, $\lambda$ acts as a knob that interpolates between \textbf{pure control design} and
\textbf{system-identification-driven design}. This is conceptually useful because it shows that direct
and indirect methods are not two completely disconnected worlds, but rather two endpoints of a
continuum of regularized formulations.

\subsubsection{Robustness-promoting regularization for trajectory parameterization}

The previous regularizer promotes equivalence with least-squares identification, but it does not
explicitly address robustness with respect to noise. To understand how robustness enters, it is
useful to go back to the noisy data relation
\[
A+BK = (X_1-W_0)G.
\]
In the certainty-equivalent formulation, one replaces this by the nominal approximation
\[
A+BK \approx X_1G,
\]
thereby neglecting the noise term $W_0G$. The issue is that this noise term is multiplied by
$G$, so large values of $G$ amplify the effect of data perturbations on the reconstructed
closed-loop matrix.\\
The same phenomenon appears in the Lyapunov inequality. The nominal direct constraint is
\[
X_1GPG^\top X_1^\top - P + I \preceq 0,
\]
whereas the true noisy relation would involve
\[
(X_1-W_0)GPG^\top (X_1-W_0)^\top - P + I \preceq 0.
\]
So the discrepancy between the nominal and true constraints depends on the matrix
\[
GPG^\top.
\]
This suggests the following principle: \textbf{for robustness, the quantity $GPG^\top$ should be
kept small}. Intuitively, if $GPG^\top$ is small, then the effect of the unknown term $W_0$ on
the closed-loop Lyapunov inequality is less pronounced.\\
This motivates the robustness-promoting regularized problem
\[
\begin{aligned}
\min_{P \succeq I,\;K,\;G}\quad
& \text{Tr}(QP)+\text{Tr}(K^\top RKP)+\lambda \|GPG^\top\| \\
\text{s.t.}\quad
& X_1GPG^\top X_1^\top - P + I \preceq 0,\\
& \begin{bmatrix}
K\\[0.3em]
I
\end{bmatrix}
=
\begin{bmatrix}
U_0\\[0.3em]
X_0
\end{bmatrix}G.
\end{aligned}
\]
The regularization term penalizes solutions for which the policy parameterization is too
sensitive to data perturbations. In contrast with the certainty-equivalence-promoting regularizer,
which selects the minimum-norm representative in the affine family of feasible $G$, this new
term explicitly aims at reducing the amplification of noise in the Lyapunov constraint.\\

There is also a useful connection between the two types of regularization. Since the nullspace
directions of $G$ are precisely those that are invisible in the consistency constraint
\[
\begin{bmatrix}
K\\[0.3em]
I
\end{bmatrix}
=
\begin{bmatrix}
U_0\\[0.3em]
X_0
\end{bmatrix}G,
\]
they may still enlarge $GPG^\top$ without changing the nominal gain $K$. For this reason, the
projection regularizer $\|\Pi G\|$ can also indirectly help to keep $GPG^\top$ small. So the
certainty-equivalence-promoting and robustness-promoting regularizers are different, but not
completely unrelated.

\subsubsection{Summary and limitations of direct trajectory parameterization}

We can now summarize the direct LQR formulation based on trajectory parameterization. The
starting point is the fact that, under persistence of excitation,
\[
\forall K\;\; \exists G \;\; \text{such that}\;\;
\begin{bmatrix}
K\\[0.3em]
I
\end{bmatrix}
=
\begin{bmatrix}
U_0\\[0.3em]
X_0
\end{bmatrix}G,
\]
and hence
\[
A+BK = (X_1-W_0)G.
\]
Neglecting the noise yields the nominal direct certainty-equivalent relation
\[
A+BK = X_1G.
\]
This leads to a direct data-driven LQR formulation in which the controller is parameterized by
the matrix $G$ instead of by an explicit model $(A,B)$.

Because $G$ is not unique, regularization is essential. Two important choices are:
\begin{itemize}
    \item the \textbf{certainty-equivalence-promoting regularizer} $\|\Pi G\|$, which selects the
    least-squares-compatible representative;
    \item the \textbf{robustness-promoting regularizer} $\|GPG^\top\|$, which tries to reduce sensitivity
    to noise.
\end{itemize}
This framework is attractive because it bypasses explicit model identification and keeps the
entire design directly at the data level. However, it also has important drawbacks.\\

The first drawback is that the \textbf{dimension of the optimization variable $G$ grows with the amount
of data}. If the experiment contains many samples, then $G$ has many rows, and the
optimization problem becomes increasingly large. So, unlike model-based or indirect approaches,
this parameterization does not have a complexity that is determined only by the state and input
dimensions.\\
The second drawback is that \textbf{regularization is not optional}. Because of the non-uniqueness
of $G$, some regularizing principle must be introduced in order to select a meaningful solution.
This is not merely a numerical issue: different regularizers correspond to different control
design philosophies, such as certainty-equivalence or robustness.\\
These observations motivate the search for alternative direct parameterizations that retain the
advantages of data-driven design but avoid the explicit dependence on the full data length. This
is exactly the role of the \emph{sample covariance parameterization}, which we discuss next.

\subsection{Direct data-driven LQR via sample covariance parameterization}

The trajectory-based parameterization uses the full data matrices $X_0$, $U_0$, and $X_1$.
A different possibility is to compress the information contained in the data through sample
covariances. This leads to a direct formulation whose dimension no longer grows with the
number of collected samples.\\

Starting from the same data equation
\[
X_1 = AX_0 + BU_0 + W_0,
\]
define the sample covariance matrices
\[
\Phi :=
\frac{1}{T}
\begin{bmatrix}
U_0\\[0.3em]
X_0
\end{bmatrix}
\begin{bmatrix}
U_0\\[0.3em]
X_0
\end{bmatrix}^\top,
\qquad
\Phi^+ :=
\frac{1}{T}
X_1
\begin{bmatrix}
U_0\\[0.3em]
X_0
\end{bmatrix}^\top.
\]
Under the persistence of excitation condition,
\[
\rank\!\begin{bmatrix}
U_0\\[0.3em]
X_0
\end{bmatrix}=n+m,
\]
the matrix $\Phi$ is positive definite:
\[
\Phi \succ 0.
\]
This is the covariance analogue of the full-row-rank property used in the trajectory
parameterization.\\
The covariance parameterization is based on the observation that, for every gain $K$, there
exists a matrix $V$ such that
\begin{equation}
\begin{bmatrix}
K\\[0.3em]
I
\end{bmatrix}
=
\frac{1}{T}
\begin{bmatrix}
U_0\\[0.3em]
X_0
\end{bmatrix}
\begin{bmatrix}
U_0\\[0.3em]
X_0
\end{bmatrix}^\top V
=
\Phi V.
\label{eq:cov-param}
\end{equation}
So the role played previously by $G$ is now taken by a matrix $V$ living in the smaller space
associated with the sample covariance matrix.\\
Using \eqref{eq:cov-param}, the closed-loop matrix becomes
\[
A+BK
=
[B\;\;A]
\begin{bmatrix}
K\\[0.3em]
I
\end{bmatrix}
=
[B\;\;A]\Phi V.
\]
From the data equation, this can be rewritten as
\[
A+BK
=
\frac{1}{T}(X_1-W_0)
\begin{bmatrix}
U_0\\[0.3em]
X_0
\end{bmatrix}^\top V.
\]
If noise is neglected, one obtains the certainty-equivalent approximation
\[
A+BK = \Phi^+ V.
\]
So, exactly as in the trajectory case, the unknown closed-loop matrix is replaced by a quantity
constructed directly from the data, but now only through the empirical covariances $\Phi$ and
$\Phi^+$.\\
Substituting this relation into the LQR parameterization gives the covariance-based direct
formulation
\[
\begin{aligned}
\min_{K,\;P \succeq I,\;V}\quad
& \text{Tr}(QP)+\text{Tr}(K^\top RKP)\\
\text{s.t.}\quad
& P = I + (\Phi^+V)P(\Phi^+V)^\top,\\
& \begin{bmatrix}
K\\[0.3em]
I
\end{bmatrix}
= \Phi V.
\end{aligned}
\]
Equivalently, since $\Phi \succ 0$, the second constraint determines $V$ uniquely as
\[
V = \Phi^{-1}
\begin{bmatrix}
K\\[0.3em]
I
\end{bmatrix}.
\]
This is a major difference with respect to the trajectory parameterization: in the covariance
parameterization the auxiliary variable is \textbf{unique}, so the nullspace ambiguity disappears.\\
This has several important consequences:
\begin{itemize}
    \item the dimension of the parameterization is \textbf{independent of the data length} $T$;
    \item the covariance matrices admit \textbf{recursive rank-one updates}, so the parameterization is
    naturally suitable for online data accumulation;
    \item the auxiliary variable $V$ is uniquely determined, so no nullspace regularization is
    needed to select among infinitely many representatives.
\end{itemize}
Moreover, this covariance formulation is again equivalent to indirect certainty-equivalent LQR.
Indeed, from
\[
\Phi^+ \Phi^{-1}
=
\frac{1}{T}
X_1
\begin{bmatrix}
U_0\\[0.3em]
X_0
\end{bmatrix}^\top
\left(
\frac{1}{T}
\begin{bmatrix}
U_0\\[0.3em]
X_0
\end{bmatrix}
\begin{bmatrix}
U_0\\[0.3em]
X_0
\end{bmatrix}^\top
\right)^{-1},
\]
one recovers the same least-squares estimate as before:
\[
[\hat B\;\;\hat A] = \Phi^+\Phi^{-1}.
\]
Therefore,
\[
\Phi^+V
=
\Phi^+\Phi^{-1}
\begin{bmatrix}
K\\[0.3em]
I
\end{bmatrix}
=
[\hat B\;\;\hat A]
\begin{bmatrix}
K\\[0.3em]
I
\end{bmatrix}
=
\hat A+\hat B K.
\]
So the covariance-based direct parameterization is another way of writing the same
certainty-equivalent closed-loop matrix, but in a compressed form that depends only on sample
covariances.

\subsubsection{Robustness-promoting regularization for covariance parameterization}

The covariance formulation removes the nullspace ambiguity, but it does not remove the effect
of noise. Indeed, the exact noisy relation is
\[
A+BK=
\frac{1}{T}(X_1-W_0)
\begin{bmatrix}
U_0\\[0.3em]
X_0
\end{bmatrix}^\top V.
\]
When this expression is substituted into the Lyapunov inequality, one obtains terms of the form
\[
\frac{1}{T^2}
(X_1-W_0)
\begin{bmatrix}
U_0\\[0.3em]
X_0
\end{bmatrix}^\top
VPV^\top
\begin{bmatrix}
U_0\\[0.3em]
X_0
\end{bmatrix}
(X_1-W_0)^\top.
\]
The nominal certainty-equivalent approximation replaces this by
\[
\frac{1}{T^2}
X_1
\begin{bmatrix}
U_0\\[0.3em]
X_0
\end{bmatrix}^\top
VPV^\top
\begin{bmatrix}
U_0\\[0.3em]
X_0
\end{bmatrix}
X_1^\top.
\]
The gap between the noisy and nominal terms is governed by the middle factor
\[
\frac{1}{T}
\begin{bmatrix}
U_0\\[0.3em]
X_0
\end{bmatrix}^\top
VPV^\top
\begin{bmatrix}
U_0\\[0.3em]
X_0
\end{bmatrix}.
\]
Its trace can be written as
\[
\text{Tr}\!\left(
\frac{1}{T}
\begin{bmatrix}
U_0\\[0.3em]
X_0
\end{bmatrix}^\top
VPV^\top
\begin{bmatrix}
U_0\\[0.3em]
X_0
\end{bmatrix}
\right)
=
\text{Tr}(\Phi V P V^\top).
\]
This suggests that, in the covariance parameterization, a natural robustness-promoting quantity
to keep small is
\[
\text{Tr}(\Phi V P V^\top).
\]
So the covariance-based counterpart of robustness regularization penalizes the energy of the
middle covariance-weighted factor through a term such as
\[
\lambda\,\text{Tr}(\Phi V P V^\top).
\]
The intuition is exactly the same as before: the smaller this quantity is, the less sensitive the
Lyapunov inequality becomes to the noise entering through $W_0$.\\

This concludes the main constructions based on direct data-driven parameterizations. We now
have two complementary viewpoints:
\begin{itemize}
    \item the \textbf{trajectory parameterization}, which is very direct and transparent but whose
    dimension grows with the data length and requires explicit regularization;
    \item the \textbf{sample covariance parameterization}, which compresses the data into fixed-size
    matrices, yields a unique auxiliary variable, and recovers the same certainty-equivalent
    solution in a more compact form.
\end{itemize}


\subsection{Comparison of data-driven LQR methods}

We conclude this part of the chapter by comparing the main data-driven LQR approaches that we
have introduced so far. The goal of this comparison is not to declare a universally best method,
but rather to highlight the trade-offs between \emph{modularity}, \emph{optimality},
\emph{robustness}, and \emph{computational size}.\\
The methods we compare are:
\begin{itemize}
    \item \textbf{indirect certainty-equivalent LQR}, where one first identifies $(\hat A,\hat B)$ and then
    solves the corresponding model-based LQR problem;
    \item \textbf{direct LQR with trajectory parameterization}, in three variants:
    without regularization, with certainty-equivalence-promoting regularization, and with
    robustness-promoting regularization;
    \item \textbf{direct LQR with sample covariance parameterization}, with and without robustness
    regularization.
\end{itemize}
A compact summary is reported in Table~\ref{tab:dd-lqr-comparison}.

\begin{table}[h!]
\centering
\small
\renewcommand{\arraystretch}{1.2}
\setlength{\tabcolsep}{4pt}
\resizebox{\textwidth}{!}{%
\begin{tabular}{|l|c|ccc|cc|}
\hline
& \textbf{Indirect CE} & \multicolumn{3}{c|}{\textbf{Direct LQR + trajectory parameterization}} & \multicolumn{2}{c|}{\textbf{Direct LQR + covariance parameterization}} \\
\cline{3-7}
& \textbf{LQR} & \textbf{No reg.} & \textbf{CE reg.} & \textbf{Robust reg.} & \textbf{No reg.} & \textbf{Robust reg.} \\
\hline
\textbf{Modular} & Yes & No & No & No & No & No \\
\hline
\textbf{Optimality} & Good & Bad & Good & Bad & Good & Good \\
\hline
\textbf{Robust stability} & Moderate & Bad & Moderate & Good & Moderate & Good \\
\hline
\textbf{Problem size} & Constant & \multicolumn{3}{c|}{Scales linearly with data size} & \multicolumn{2}{c|}{Constant} \\
\hline
\end{tabular}%
}
\caption{Comparison of the main data-driven LQR methods discussed in this chapter.}
\label{tab:dd-lqr-comparison}
\end{table}
We now explain the meaning of each row in the table.

\paragraph{Modularity.}
The indirect certainty-equivalent approach is the only \textbf{modular} one. Indeed, it keeps the
classical separation
\[
\text{identification} \;\longrightarrow\; \text{control design},
\]
so the modeling step and the controller synthesis step can be developed and analyzed
independently. By contrast, the direct methods are \textbf{non-modular}: data and control design are
intertwined in a single optimization problem. This lack of modularity is not necessarily a
drawback in itself, but it does make the overall design less transparent.

\paragraph{Optimality.}
The optimality row refers to how well the resulting controller aligns with the
certainty-equivalent LQR objective.\\

The \textbf{indirect CE-LQR} method is marked as \emph{good} because it computes exactly the LQR
controller associated with the least-squares identified model. So, within the certainty-equivalent
philosophy, it is the natural benchmark.\\
For \textbf{trajectory parameterization without regularization}, the performance is marked as
\emph{bad}. The reason is that the auxiliary variable $G$ is not unique, and without any principle
to select among the feasible solutions, the optimization may exploit unfavorable nullspace
directions. This generally breaks the connection with least-squares identification and leads to a
controller that is no longer well aligned with the nominal LQR objective.\\
When \textbf{certainty-equivalence-promoting regularization} is added, optimality becomes
\emph{good}, because the projector regularizer $\|\Pi G\|$ drives the solution toward the
minimum-norm representative associated with least-squares identification. Hence the direct
trajectory formulation recovers the same nominal solution as the indirect approach.\\
With \textbf{robustness-promoting regularization} in the trajectory formulation, optimality is again
marked as \emph{bad}. This does not mean the controller is useless; rather, it means that the
design is no longer trying to reproduce the certainty-equivalent optimum. Instead, it sacrifices
nominal optimality in order to improve robustness.\\
For the \textbf{covariance parameterization}, optimality is marked as \emph{good} both without and
with robustness regularization. Without regularization, this follows from the fact that the
covariance parameterization is uniquely determined and is equivalent to indirect
certainty-equivalent LQR. With robustness regularization, the table still marks the method as
good, reflecting the fact that the covariance formulation preserves a much cleaner and more
well-posed connection with the certainty-equivalent solution than the unregularized trajectory
formulation.

\paragraph{Robust stability.}
The robustness row summarizes how reliable the different methods are in the presence of noisy
data and modeling errors.\\

For \textbf{indirect CE-LQR}, robustness is classified as \emph{moderate}. This is consistent with the
discussion above: certainty-equivalence is simple and often effective, but it does not explicitly
optimize robustness unless additional uncertainty-aware design steps are introduced.\\
For \textbf{trajectory parameterization without regularization}, robustness is \emph{bad}. This is the
least protected formulation, since the non-uniqueness of $G$ may amplify noise significantly and
there is no mechanism preventing harmful solutions.\\
With \textbf{certainty-equivalence-promoting regularization}, robustness improves to
\emph{moderate}. In this case, the direct method essentially reproduces the least-squares indirect
solution, so its robustness properties are similar to those of indirect CE-LQR.\\
With \textbf{robustness-promoting regularization}, the trajectory-based direct method becomes
\emph{good} from the robustness viewpoint, because the added penalty explicitly aims at reducing
quantities such as $GPG^\top$, which control the sensitivity of the Lyapunov constraint to noise.\\
The same logic applies to the \textbf{covariance parameterization}. Without extra regularization, its
robustness is \emph{moderate}, since it is again tied to certainty-equivalence. Once a suitable
robustness-promoting penalty is added, robustness becomes \emph{good}.

\paragraph{Problem size.}
The last row concerns the computational dimension of the optimization problem.\\

For \textbf{indirect CE-LQR}, the problem size is \textbf{constant} once the model has been identified,
because the controller synthesis depends only on the state and input dimensions $(n,m)$.\\
For the \textbf{trajectory-parameterized direct methods}, the problem size \textbf{scales linearly with the
amount of data}. This is because the matrix $G$ has as many rows as there are data points in the
experiment. Hence, as more samples are collected, the optimization problem grows.\\
For the \textbf{covariance-parameterized direct methods}, the problem size is again \textbf{constant},
because the data enter only through the fixed-size matrices $\Phi$ and $\Phi^+$. This is one of the
main reasons why the covariance formulation is attractive in practice.\\

The table makes the trade-offs quite visible.\\
Indirect certainty-equivalent LQR is attractive because it is modular, has fixed problem size,
and provides a good nominal solution, but its robustness is only moderate.\\
Direct trajectory parameterization is conceptually very flexible and can be made robust, but it
has two main weaknesses: its dimension grows with data size, and regularization is essential.\\
Direct covariance parameterization keeps the direct data-driven spirit while avoiding the growth
with data size. It also preserves the connection with certainty-equivalent design in a much more
compact way, and it can be enhanced with robustness-promoting regularization.\\

So, roughly speaking:
\begin{itemize}
    \item \textbf{indirect CE-LQR} is the simplest modular baseline;
    \item \textbf{direct trajectory parameterization} offers the most design freedom, but at a higher
    computational and regularization cost;
    \item \textbf{direct covariance parameterization} provides a compact compromise between directness,
    nominal performance, and tractability.
\end{itemize}
This comparison will serve as a useful reference point for interpreting the strengths and
limitations of the different data-driven LQR constructions developed in this chapter.